\expandafter\gdef\csname SNside1\endcsname{\node[anchor=north west,inner sep=0pt,text width=48mm,font=\SideFont,align=left] at ([xshift=138mm,yshift=-516.15pt]current page.north west) {\textcolor{SNink}{\textbf{课程连接 · 功能部件接到数据通路}}\par\vspace{1.5mm}把 c=a+b 分成取a、取b、相加、存c后，第5章会逐步标出寄存器传送和控制信号。这里的功能划分仍成立；一块芯片可以同时实现多种功能，不能按芯片数倒推五部件。};}
\expandafter\gdef\csname SNside2\endcsname{\node[anchor=north west,inner sep=0pt,text width=48mm,font=\SideFont,align=left] at ([xshift=138mm,yshift=-237.14pt]current page.north west) {\textcolor{SNink}{\textbf{最小反例 · 同ISA不保证文件通用}}\par\vspace{1.5mm}最小反例：两系统使用同一指令集，但系统调用号、动态库和可执行文件格式不同。一份应用文件仍可能无法直接运行。ISA兼容解决指令语义，运行环境兼容还需操作系统与ABI条件。};}
\expandafter\gdef\csname SNside3\endcsname{\node[anchor=north west,inner sep=0pt,text width=48mm,font=\SideFont,align=left] at ([xshift=138mm,yshift=-95.00pt]current page.north west) {\textcolor{SNink}{\textbf{不变量 · 地址单位必须统一}}\par\vspace{1.5mm}顺序执行的后继地址应等于当前地址加指令占用的地址单位数。按字节编址且指令4B时加4；若题设每次只取1B，PC可逐次加1。先统一单位，才不会把局部取指步骤误作整条指令长度。};}
\expandafter\gdef\csname SNside4\endcsname{\node[anchor=north west,inner sep=0pt,text width=48mm,font=\SideFont,align=left] at ([xshift=138mm,yshift=-95.00pt]current page.north west) {\textcolor{SNink}{\textbf{错项排除 · 汇编符号不能直接执行}}\par\vspace{1.5mm}CPU读取的是按ISA编码的机器位串。汇编语言中的助记符还需编码，宏与伪指令可能展开为多条或零条机器指令。题目问``直接执行''时，应排除仅便于人阅读的符号表示。};}
\expandafter\gdef\csname SNside5\endcsname{\node[anchor=north west,inner sep=0pt,text width=48mm,font=\SideFont,align=left] at ([xshift=138mm,yshift=-95.00pt]current page.north west) {\textcolor{SNink}{\textbf{最小反例 · 速度翻倍不等于总时减半}}\par\vspace{1.5mm}最小反例：10秒任务中只有2秒可加速，即使这2秒缩短为0，总时间仍至少8秒。只有整个统计对象都按相同倍率加速，总历时才按该倍率缩短；先画出未改变的部分。\par\smallskip\begin{tikzpicture}[x=1mm,y=1mm,>=Stealth,node distance=5mm]
\node[draw=SNink!50,fill=SNblue,rounded corners=1mm,text width=34mm,align=center,inner sep=1.2mm,font=\fontsize{7.5}{9}\selectfont] (m0) at (0,0) {原10秒};
\node[draw=SNink!50,fill=SNblue,rounded corners=1mm,text width=34mm,align=center,inner sep=1.2mm,font=\fontsize{7.5}{9}\selectfont] (m1) at (0,-11) {8秒不变};
\node[draw=SNink!50,fill=SNblue,rounded corners=1mm,text width=34mm,align=center,inner sep=1.2mm,font=\fontsize{7.5}{9}\selectfont] (m2) at (0,-22) {至少8秒};
\draw[->,SNink] (m0) -- (m1);
\draw[->,SNink] (m0) edge[bend left=40] (m2);
\end{tikzpicture}
\par{\fontsize{7.5}{9}\selectfont 可加速部分即使归零，总时间仍不小于8秒。}};}
\expandafter\gdef\csname SNside6\endcsname{\node[anchor=north west,inner sep=0pt,text width=48mm,font=\SideFont,align=left] at ([xshift=138mm,yshift=-95.00pt]current page.north west) {\textcolor{SNink}{\textbf{课程连接 · Cache和操作系统怎样进入时间式}}\par\vspace{1.5mm}Cache缺失与流水线阻塞通常进入平均CPI；编译优化会改变动态指令数和组成；调度与I/O等待可能进入总历时。连接这些课程时先确认统计口径，避免把所有等待都塞入CPU执行时间。};}
\expandafter\gdef\csname SNside7\endcsname{\node[anchor=north west,inner sep=0pt,text width=48mm,font=\SideFont,align=left] at ([xshift=138mm,yshift=-95.00pt]current page.north west) {\textcolor{SNink}{\textbf{最小反例 · 字长不能推出地址容量}}\par\vspace{1.5mm}同样是32位整数处理器，可以配较小物理内存，也可以有更宽物理地址。最小反例已表明``字长32位所以最多4GiB''缺少地址位数前提。容量必须由地址数与每地址对应的单位共同确定。};}
\expandafter\gdef\csname SNside8\endcsname{\node[anchor=north west,inner sep=0pt,text width=48mm,font=\SideFont,align=left] at ([xshift=138mm,yshift=-95.00pt]current page.north west) {\textcolor{SNink}{\textbf{不变量 · 平均值守住总量}}\par\vspace{1.5mm}无论指令如何分类，平均CPI乘总动态指令数都应还原总周期数。若换一种分组后总周期发生变化，说明权重、分类或计数口径出了问题；这个回代比死记加权公式更容易排错。};}
\AddToShipoutPictureFG{\begin{tikzpicture}[remember picture,overlay]\ifcsname SNside\thepage\endcsname\csname SNside\thepage\endcsname\fi\end{tikzpicture}}